%! Choosing a Sample: Four Sampling Techniques — Unit 1, Lesson 1.3 — Guided Notes & Practice (Answer Key)
% PILOT SHAPE (2026-09-06): modeled on AP Statistics lesson 1.4 minus its AP Practice page.
% THE in-class document. The notes carry the whole gradual release, closed by a SPOKEN debrief.
% They run SHORT on purpose: the period ends with students starting the homework, which is
% where the individual practice lives.
%   I DO   : the vocabulary box and notes sections 1-2.                              (~20 min)
%   WE DO  : the practicebox ("Guided Practice") — Cedar Ridge Apartments, together.  (~15 min)
%   YOU DO : nothing on this page — ../homework/ IS the individual practice, started in the
%            last ten minutes of the period.
% PAGE PLAN: p1 = vocab box + hook (no objective box — the cover carries the targets);
% p2 = section 1; p3 = section 2; p4 = the Guided Practice alone. Every fill-in cluster opens
% with a \wordbank{}. No closing practice set, no extension.
% Vocabulary is GIVEN here, not discovered. THE CRUX is the second half of section 2: "chance
% chose" is not enough — a cluster sample only works when each cluster is a small mix of the
% whole, and a Riverbend homeroom is alike inside (a stratum), so Maya's two random homerooms
% are 60 ninth graders and zero seniors (PS.DC.2a-b). The second half of section 1 carries the
% second trap: "random" does not mean "unplanned" — the test is WHO COULD NEVER HAVE BEEN PICKED.
%
% THE DATA
%   Riverbend HS: 900 students, the council surveys 60. Grades 270/240/210/180 -> at 1/15:
%     18/16/14/12. Digit strip 7 1 4 0 2 8 9 5 1 3 6 6 -> students 714, 28, (skip 951), 366.
%     Systematic k = 900/60 = 15, start 7 -> 7, 22, 37, 52 (the warm-up's numbers).
%     30 homerooms of 30: homerooms 1-9 hold the 9th grade, 10-17 the 10th, 18-24 the 11th,
%     25-30 the 12th. Maya draws homerooms 4 and 7 -> 2 x 30 = 60 students, all 9th graders.
%     About one draw in four (200/870 = 23%) lands both rooms in a single grade.
%   Guided Practice, Cedar Ridge Apartments: 600 households in 30 buildings of 20; buildings
%     1-6 studios (120), 7-20 one-bedrooms (280), 21-30 two-bedrooms (200); 60 households at
%     1/10 -> 12/28/20; a cluster draw of buildings 2, 5, 6 -> 60 studio households.
%
% PAGE LOCKSTEP with notes_key/: the key is GENERATED from this file — every \blank{} becomes
% an \ans{}, every \par\writelines{n} becomes n \ansline{} calls, every \vterm{} a
% \vtermans{}{}. Every \begin{work} block is byte-identical. Blank widths are sized to their
% answers so the two files wrap alike.
\documentclass[12pt]{article}
\usepackage{saar-article}
\usepackage{saar-key}   % pulls in -boxes; do NOT also load -boxes
\newcommand{\TallMath}[1]{$\displaystyle #1\rule[-1.4em]{0pt}{3.2em}$}
% Word bank strip for every fill-in cluster (user request 2026-09-06). Defined per document;
% the key inherits it and prints the same strip, so heights match.
\newcommand{\wordbank}[1]{\par\vspace{2pt}\noindent\fcolorbox{goldacc}{goldbg}{\parbox{\dimexpr\linewidth-2\fboxsep-2\fboxrule\relax}{\small\textbf{Word bank:} #1}}\par\vspace{4pt}}
% Vocabulary rows: a FIXED-HEIGHT row carrying the term label and OPEN writing space —
% no underline beside the term and no rule line (user correction 2026-09-06: the black
% answer line crowded the row; students write beside and below the term). The key fills
% exactly the same height, so the box cannot drift blank vs. keyed. saar-article's
% \termblank adds an inline blank and a rule line, so the pair is defined here (the key is
% generated from this file and inherits both). 2.3cm per row gives students three
% handwritten lines of room; keep key definitions to two lines.
\newcommand{\termrowheight}{2.3cm}
\newlength{\termrowinset}\setlength{\termrowinset}{7pt}
\newcommand{\vterm}[1]{%
  \par\noindent\begin{minipage}[t][\termrowheight][t]{\linewidth}%
    \vspace*{\termrowinset}%
    \textbf{\textcolor{cerulean}{#1:}}%
  \end{minipage}\par}
\newcommand{\vtermans}[2]{%
  \par\noindent\begin{minipage}[t][\termrowheight][t]{\linewidth}%
    \vspace*{\termrowinset}%
    \textbf{\textcolor{cerulean}{#1:}}\quad\ans{#2}%
  \end{minipage}\par}

\begin{document}

\pageheader{Unit 1, Lesson 1.3}{Guided Notes \& Practice --- Answer Key}
% NAMESTRIP: no \namedateperiod here — mirrors the blank. NO teachernote here —
% teacher prose lives in the lesson plan.

\vspace{0.15in}

\begin{vocabbox}
\small
Fill in each term \textbf{as we name it} in the notes below.
\vtermans{Sampling frame}{The list of individuals a study can actually choose from. Anyone left off the frame can never be sampled, however good the method.}
\vtermans{Simple random sample (SRS)}{Every individual --- and every possible group of $n$ individuals --- has an equal chance of being chosen. Number the frame and draw, skipping out-of-range and repeat values.}
\vtermans{Stratified sample}{Split the population into strata (alike inside, different between); take a separate random sample from every stratum, in proportion.}
\vtermans{Systematic sample}{From an ordered list, take every $k$th individual after a random start, where $k = N \div n$. Only the start is random.}
\vtermans{Cluster sample}{Draw a few whole groups (clusters) at random and collect data from everyone in them. Works only when each cluster is a small mix of the whole.}
\end{vocabbox}

\boxguard[12]
\begin{hookbox}
\small
The student council needs opinions from $60$ of Riverbend's $900$ students about the late bus.
Maya puts the $30$ homeroom numbers in a hat, draws two --- homerooms $4$ and $7$ --- and
surveys every one of the $30$ students in each room. Pure chance picked those rooms.

\vspace{3pt}
\textbf{Is Maya's sample as trustworthy as drawing $60$ student names out of the hat?} Circle
one: \quad yes \qquad no \qquad it depends \quad --- and say why you think so. We vote again
at the end of section~2.
\par\ansline{Answers vary --- most of the room says yes (``it was random'') or it depends.}
\ansline{(Settled at the end of section 2: no --- both rooms are 9th grade; not one senior was asked.)}
\end{hookbox}

% ── I DO: section 1 ──────────────────────────────────────────────────────────
% Page 2: section 1 (page 1 is the vocab box + hook).
\newpage
\begin{notesbox}{1. Who Chooses the Sample --- and How to Let Chance Do It}
\small
Every study so far had a sample. None of them said \emph{how} it was picked. That decision is
made before any data exist, and it decides what the statistic will be worth.
\wordbank{chance \;\textbullet\; convenience \;\textbullet\; could not \;\textbullet\; yes \;\textbullet\; no --- skip it \;\textbullet\; student 28 \;\textbullet\; student 366 \;\textbullet\; none \;\textbullet\; three \;\textbullet\; start \;\textbullet\; 15 \;\textbullet\; 22 \;\textbullet\; 37 \;\textbullet\; 52}
The \textbf{sampling frame} is the list of individuals a study can actually choose from --- at
Riverbend, the office's list of all $900$ students. Once you have a frame there are two ways to
choose. In a \textbf{probability sample}, \ans{chance} does the choosing: every individual on
the frame has a known chance of being picked. In a \ans{convenience} sample, the researcher takes
whoever is easiest to reach.

\vspace{3pt}
Who chose each sample --- chance, or convenience?
\nopagebreak

\begin{center}
\renewcommand{\arraystretch}{1.35}
\begin{tabularx}{\linewidth}{@{} X p{2.9cm} @{}}
\toprule
\textbf{How the sample was taken} & \textbf{Who chose it?} \\
\midrule
A pool staff member asks the $75$ members who arrive on Tuesday evening. & \ans{convenience} \\
All $900$ students are numbered and a computer picks $60$ numbers. & \ans{chance} \\
A teacher surveys her own third-period class. & \ans{convenience} \\
The $30$ homeroom numbers go in a hat and $2$ are drawn. & \ans{chance} \\
\bottomrule
\end{tabularx}
\end{center}

\vspace{3pt}
\begin{center}
\fcolorbox{redacc}{redbg}{\parbox{0.94\linewidth}{\small
\textbf{Careful --- \emph{random} does not mean \emph{unplanned}.} The teacher did not pick
favorites, but a fourth-period student \ans{could not} have been chosen, so chance did not do
the choosing. The test for every sample: \emph{who could never have been picked?} If that list
is not empty, it is a \ans{convenience} sample --- and no size fixes it.}}
\end{center}

\vspace{6pt}
\textbf{\textcolor{cerulean}{Simple Random Sample --- and the Systematic Shortcut.}}
In a \textbf{simple random sample} of size $n$, every individual --- and every possible group
of $n$ individuals --- has an equal chance of being chosen. \textbf{How:} number the frame $1$
to $N$, run a random number generator, and read numbers off, \textbf{skipping} any that are out
of range or already used, until you have $n$ different individuals.

\vspace{3pt}
Riverbend's students are numbered $001$--$900$. The generator returns
{\ttfamily 7\,1\,4\;\;0\,2\,8\;\;9\,5\,1\;\;3\,6\,6}. Read the digits three at a time and decide each group.
\nopagebreak

\begin{center}
\renewcommand{\arraystretch}{1.35}
\begin{tabularx}{\linewidth}{@{} c p{3.6cm} X @{}}
\toprule
\textbf{Digits} & \textbf{In range $001$--$900$?} & \textbf{Student chosen} \\
\midrule
714 & yes & student 714 \\
028 & \ans{yes} & \ans{student 28} \\
951 & \ans{no --- skip it} & \ans{none} \\
366 & \ans{yes} & \ans{student 366} \\
\bottomrule
\end{tabularx}
\end{center}

\vspace{3pt}
Why read three digits at a time? Because the largest number on the frame, $900$, has
\ans{three} digits. An SRS is the fairest of the four and the one the others are measured
against. Its price: you need a complete list, and the $60$ it picks are scattered across the
building.

\vspace{3pt}
If the frame is already an ordered list, a \textbf{systematic sample} skips the $60$ separate
draws: compute the interval $k = N \div n$, pick a random \ans{start} between $1$ and $k$,
then take every $k$th individual. For Riverbend, $k = 900 \div 60 =$ \ans{15} (the
warm-up), the start lands on $7$, and the picks are $7$, \ans{22}, \ans{37},
\ans{52}, \dots\ Only the \textbf{start} is random --- after that the list decides, which is
fine as long as the order of the list has nothing to do with the interval.

\end{notesbox}

% Page 3: section 2.
\newpage
\begin{notesbox}{2. Sampling in Groups --- Some From Every Group, or Everyone From a Few}
\small
Most populations come in groups: grades, homerooms, buildings. Two techniques use the groups
--- the two everybody mixes up, because \emph{which one you may use depends on how the groups
are built}.
\wordbank{every \;\textbullet\; all \;\textbullet\; a few \;\textbullet\; everyone \;\textbullet\; a random sample \;\textbullet\; mix \;\textbullet\; alike \;\textbullet\; stratum \;\textbullet\; stratified \;\textbullet\; 9th \;\textbullet\; 0 \;\textbullet\; Yes \;\textbullet\; No \;\textbullet\; 16 \;\textbullet\; 14 \;\textbullet\; 12 \;\textbullet\; 60}
Groups whose members are \textbf{alike inside} and \textbf{different from group to group} are
called \textbf{strata} (one is a \emph{stratum}). In a \textbf{stratified sample} you take a
separate random sample from \ans{every} stratum, sized so each keeps its share of the
population. Seniors drive, so grade is a natural set of strata for a late-bus question. The
council wants $60$ of $900$, so take $60/900 = 1/15$ of each grade.
\nopagebreak

\begin{center}
\renewcommand{\arraystretch}{1.2}
\begin{tabularx}{\linewidth}{@{} l c X @{}}
\toprule
\textbf{Grade} & \textbf{Students} & \textbf{Number to survey} \\
\midrule
9th   & 270 & $270 \div 15 = 18$ \\
10th  & 240 & \ans{$240 \div 15 = 16$} \\
11th  & 210 & \ans{$210 \div 15 = 14$} \\
12th  & 180 & \ans{$180 \div 15 = 12$} \\
\midrule
\textbf{Total} & \textbf{900} & \ans{$18+16+14+12 = 60$} \\
\bottomrule
\end{tabularx}
\end{center}

\begin{work}
  \dfrac{60}{900} &= \dfrac{1}{15} \\
  \dfrac{1}{15} \times 240 &= 16 \text{ students}
\end{work}

Stratifying \textbf{guarantees} every grade is heard, in proportion. A fair SRS only makes that
\emph{likely} --- by chance it could land on very few seniors.

\vspace{5pt}
\textbf{\textcolor{cerulean}{Cluster Sample --- and Maya's Homeroom Draw.}}
When there is no list of individuals, only a list of groups, a \textbf{cluster sample} draws
\ans{a few} whole groups at random and collects data from \ans{everyone} in them. It is the
cheapest of the four --- and it works only when each cluster is a small \ans{mix} of the
whole. Riverbend's $900$ students sit in $30$ homerooms of $30$; Maya drew rooms $4$ and $7$.

\begin{work}
  2 \times 30 &= 60 \text{ students}
\end{work}

Now the fact Maya did not check: \textbf{homerooms $1$--$9$ hold the 9th grade, $10$--$17$
the 10th, $18$--$24$ the 11th, $25$--$30$ the 12th.} Rooms $4$ and $7$ are both
\ans{9th}-grade rooms, so the number of seniors in Maya's sample is \ans{0}. Did
chance choose? \ans{Yes} \quad Did anyone count wrong? \ans{No} \quad \textbf{Vote again.}

\vspace{3pt}
\begin{center}
\fcolorbox{cerulean}{frost}{\parbox{0.94\linewidth}{\small
\textbf{``Chance chose'' is not enough --- the technique has to match the groups.} A Riverbend
homeroom is \ans{alike} inside, so it is really a \ans{stratum}, and drawing a few whole
strata can miss a whole grade --- about one draw in four lands both rooms in one grade. Alike
inside means \emph{some from \ans{every} group} --- a \ans{stratified} sample. Cluster sampling
belongs to groups that are each a small mix of the whole.}}
\end{center}

\vspace{2pt}
\nopagebreak
\begin{center}
\renewcommand{\arraystretch}{1.25}
\begin{tabularx}{\linewidth}{@{} p{3.4cm} X X @{}}
\toprule
 & \textbf{Stratified} & \textbf{Cluster} \\
\midrule
The groups are\dots & alike inside, different from each other & each a small mix of the whole \\
You use\dots & \ans{all} of the groups & \ans{a few} of the groups \\
From each you take\dots & \ans{a random sample} & \ans{everyone} \\
\bottomrule
\end{tabularx}
\end{center}

\end{notesbox}

% Page 4: the Guided Practice, alone.
\newpage
\begin{practicebox}
\small
\textbf{Cedar Ridge Apartments --- we work this one together.} Cedar Ridge has $600$
households in $30$ buildings of $20$; management wants to survey $60$ about parking. Buildings
$1$--$6$ are studios ($120$ households), $7$--$20$ one-bedrooms ($280$), $21$--$30$
two-bedrooms ($200$) --- and two-bedroom households own the most cars.
\wordbank{simple random \;\textbullet\; stratified \;\textbullet\; systematic \;\textbullet\; cluster \;\textbullet\; convenience \;\textbullet\; 12 \;\textbullet\; 28 \;\textbullet\; 20 \;\textbullet\; 60 \;\textbullet\; studio \;\textbullet\; stratum}
\textbf{(a)} Name the technique in each plan.
\nopagebreak

\begin{center}
\renewcommand{\arraystretch}{1.2}
\begin{tabularx}{\linewidth}{@{} X p{2.9cm} @{}}
\toprule
\textbf{The plan} & \textbf{Technique} \\
\midrule
Number all $600$ households and let a computer draw $60$ numbers. & \ans{simple random} \\
Draw $3$ building numbers from a hat; survey all $20$ households in each. & \ans{cluster} \\
A random sample from each unit type --- studios, one-, two-bedrooms --- in proportion. & \ans{stratified} \\
Start at a random household on the resident list, then take every $10$th one. & \ans{systematic} \\
Stand by the mailboxes after work and ask the first $60$ people who walk past. & \ans{convenience} \\
\bottomrule
\end{tabularx}
\end{center}

\vspace{3pt}
\textbf{(b)} The stratified plan: $60$ of $600$ is $1/10$, so take $1/10$ of each unit type.
Show the one-bedroom row.
\nopagebreak

\begin{center}
\renewcommand{\arraystretch}{1.2}
\begin{tabularx}{\linewidth}{@{} l c X @{}}
\toprule
\textbf{Unit type} & \textbf{Households} & \textbf{Number to survey} \\
\midrule
Studio      & 120 & \ans{$120 \div 10 = 12$} \\
One-bedroom & 280 & \ans{$280 \div 10 = 28$} \\
Two-bedroom & 200 & \ans{$200 \div 10 = 20$} \\
\midrule
\textbf{Total} & \textbf{600} & \ans{$12+28+20 = 60$} \\
\bottomrule
\end{tabularx}
\end{center}

\begin{work}
  \dfrac{60}{600} &= \dfrac{1}{10} \\
  \dfrac{1}{10} \times 280 &= 28 \text{ households}
\end{work}

\vspace{3pt}
\textbf{(c)} Suppose management had run the cluster plan and the hat gave buildings $2$, $5$,
and $6$. Every one of the $60$ households surveyed is a \ans{studio}. Is a building a small
mix of the complex, or really a \ans{stratum}? Explain what the draw got wrong even though
chance chose the buildings.
\par\ansline{A stratum: each building holds one unit type, so it is alike inside, not a mix.}
\ansline{Drawing three whole strata left every one- and two-bedroom household --- the car owners --- out.}

\vspace{3pt}
\textbf{(d)} Which plan should management run, and why does it deliver what a parking survey
needs? Write the reason as \emph{``\dots because the groups are \dots''}
\par\ansline{Stratified, because the groups (unit types) are alike inside and differ on cars;}
\ansline{sampling every stratum in proportion guarantees the two-bedroom households are heard.}

\vspace{3pt}
\textbf{(e)} Management has a list of \emph{buildings} but has never kept a list of
\emph{residents}. Which of the four techniques can they still run without building that list?
\par\ansline{Only cluster --- it needs a list of groups. SRS, stratified, and systematic need every household numbered.}
\end{practicebox}

\end{document}
